\documentclass[a4paper,12pt]{article}

\input{../../Preambulo.tex}
\newcommand{\assunto}{FUNÇÕES AFINS}
\newcommand{\atividade}{atividade} % prova ou lista ou as ou ad ou atividade
\newcommand{\NumQObj}{6}
\newcommand{\ValorProva}{2}
\pgfmathparse{\ValorProva/\NumQObj}
\edef\ValorQObj{\pgfmathresult}
%\newcommand{\ValorQObj}{0.3}
%\pgfmathparse{int(round(\NumQObj * \ValorQObj))}
%\edef\ValorProva{\pgfmathresult}
\newcommand{\tipo}{dupla} % dupla ou individual
\newcommand{\calculadora}{nao} % sim ou nao
\newcommand{\turma}{ELETROELETRÔNICA 1.\textordmasculine\ ANO}
\newcommand{\numProva}{A}
\newcommand{\professor}{Igor Martins Silva}
\newcommand{\bimestre}{3}
\newcommand{\ano}{2026}
\newcommand{\mesTexto}{agosto}
\newcommand{\mesNum}{08}
\newcommand{\dia}{27}
\newcommand{\dataTexto}{\dia\ de \mesTexto\ de \ano}
\newcommand{\dataNun}{\dia/\mesNum/\ano}

\newcommand{\titulo}{}

\ifdefstring{\atividade}{prova}
  {\renewcommand{\titulo}{Avaliação de Matemática}}
{
\ifdefstring{\atividade}{lista}
  {\renewcommand{\titulo}{Lista de exercícios de Matemática}}
{
\ifdefstring{\atividade}{as}
  {\renewcommand{\titulo}{Avaliação Somativa}}
{
\ifdefstring{\atividade}{ad}
  {\renewcommand{\titulo}{Avaliação Diagnóstica}}
  {\renewcommand{\titulo}{Atividade de Matemática}}
}}}

\newcommand{\emtipo}{%
  \ifdefstring{\tipo}{individual}{\tipo}{em \tipo}%
}

\edef\pathCapa{\currfiledir}

\newcommand{\subtitulobase}[3]{%
    % #1 = título da 1ª coluna
    % #2 = conteúdo da 1ª coluna
    % #3 = mostra número da prova? (vazio ou algo)
    
    \noindent
    \begin{minipage}{2.8cm}
        \includegraphics[width=2.8cm]{\pathCapa Logo_Cefet.png}
    \end{minipage}
    \hfill
    \begin{minipage}{\dimexpr\textwidth-6.2cm\relax}
        \begin{center}
        CEFET-MG -- Campus Contagem \\[3pt]
        \titulo\ -- \bimestre.\textordmasculine\ bimestre de \ano \\[3pt]
        \professor
    \end{center}
    \end{minipage}
    \hfill
    \begin{minipage}{3.2cm}
        \includegraphics[width=3.2cm]{\pathCapa Brasao_Brasil.png}
    \end{minipage}
    
    \begin{center}
        \vspace{15pt}
        \begin{tabular}{|C{210pt}|C{70pt}|C{210pt}|}
            \hline
            \textbf{#1} &
            \textbf{DATA} &
            \textbf{TURMA} \\
            \hline
            #2 & \dataNun & \turma \\
            \hline
        \end{tabular}
    \end{center}
    
    #3
    
    \vspace{15pt}
}

\newcommand{\subtitulolis}{%
  \subtitulobase{ASSUNTO}{\assunto}{}%
}

\newcommand{\subtitulopr}{%
  \subtitulobase{ASSUNTO}{\assunto}{\boxed{\numProva}}%
}

\newcommand{\subtituloas}{%
  \subtitulobase{DISCIPLINA}{MATEMÁTICA}{\boxed{\numProva}}%
}

\newcommand{\subtituload}{%
  \subtitulobase{DISCIPLINA}{MATEMÁTICA}{}%
}

\newcommand{\valorquestao}{}

\newcommand{\definevalorquestao}{%
    \ifdefstring{\atividade}{lista}{%
        % Não faz nada para lista
    }{%
        \ifdefstring{\modo}{objetivo}{%
            \renewcommand{\valorquestao}{\ValorQObj ponto}
        }{%
            \renewcommand{\valorquestao}{\ValorQDisc pontos}
        }%
    }%
}

\newenvironment{exercicioBanco}[1][]{%
    \ifdefstring{\atividade}{lista}{%
        \begin{exercicio}%
    }{%
        \begin{exercicio}[#1]%
    }%
}{%
    \end{exercicio}
}

\ifdefstring{\atividade}{lista}{
    \pagecolor{background-color}
}{
    
}



\hypersetup{
    pdfauthor={\professor},
    pdftitle={\titulo \-- \assunto \-- \turma \-- \numProva},
}

\title{\titulo \-- \assunto \-- \turma \-- \numProva}
\author{\professor}
\date{\data}

%************* INÍCIO DO DOCUMENTO *************
\begin{document}
    \NoBgThispage
    \pagenumbering{alph}
    %
        
    \phantomsection
    \addcontentsline{toc}{section}{Capa}
    \input{../../AT_Capa_Gabarito}
    \newpage
    \null
    \thispagestyle{empty}
    \addtocounter{page}{1}
    \newpage
        
    \NoBgThispage
    \pagenumbering{arabic}
    %
        
    \subtitulopr
    
    \vspace{15pt}
    \newcommand{\modo}{objetivo} % objetivo ou discursivo
    
    \phantomsection
    \addcontentsline{toc}{section}{Exercício 1}
    \input{../FA_Exercicio2_2}
    \vspace{15pt}
    
    \phantomsection
    \addcontentsline{toc}{section}{Exercício 2}
    %%%%%%%%%%%%%%%%%%%%%%%%%%%%%%%%%
%
%       EXERCÍCIO
%
%%%%%%%%%%%%%%%%%%%%%%%%%%%%%%%%%

\ifdefstring{\atividade}{prova}{%
    \ifdefstring{\modo}{objetivo}{%
        \renewcommand{\valorquestao}{\ValorQObj\ ponto}
    }{%
        \renewcommand{\valorquestao}{\ValorQDisc\ pontos}
    }%
}%

\begin{exercicioBanco}[\valorquestao]
Determine o gráfico da função afim \(f: \bbR \to \bbR\), dada por \(f(x) = -2x + 1\).

% Define as alternativas
\newcommand{\alternativas}{%
\begin{center}
    \begin{tabular}{C{20pt}C{140pt}C{20pt}C{140pt}C{20pt}C{140pt}}
        (a) 
        &
        \begin{tikzpicture}[scale=0.5]
            % Eixos
            \draw[->] (-3.0,0) -- (3.0,0) node[right] {\(x\)};
            \draw[->] (0,-3.0) -- (0,3.0) node[above] {\(y\)};

            % Gráfico correto
            \draw[domain=-1.5:1.5, thick, blue] plot (\x, {-2*(\x) + 1});

            % Raiz
            \node[above right] at (0.5,0) {\small\(\frac{1}{2}\)};
            \node at (0.5,0) {\small\(\bullet\)};

            % Intercepto
            \node[left] at (0,1) {\small\(1\)};
            \node at (0,1) {\small\(\bullet\)};
        \end{tikzpicture}
        &
        (b) 
        &
        \begin{tikzpicture}[scale=0.5]
            % Eixos
            \draw[->] (-3.0,0) -- (3.0,0) node[right] {\(x\)};
            \draw[->] (0,-3.0) -- (0,3.0) node[above] {\(y\)};

            % Gráfico errado (inclinação positiva)
            \draw[domain=-1.5:1.5, thick, blue] plot (\x, {2*(\x) + 1});

            \node[right] at (0,1) {\small\(1\)};
            \node at (0,1) {\small\(\bullet\)};
            \node[above left] at (-0.5,0) {\small\(-\frac{1}{2}\)};
            \node at (-0.5,0) {\small\(\bullet\)};
        \end{tikzpicture}
        &
        (c)
        &
        \begin{tikzpicture}[scale=0.5]
            % Eixos
            \draw[->] (-3.0,0) -- (3.0,0) node[right] {\(x\)};
            \draw[->] (0,-3.0) -- (0,3.0) node[above] {\(y\)};

            % Erro no termo independente
            \draw[domain=-1.5:1.5, thick, blue] plot (\x, {-2*(\x) - 1});

            \node[left] at (0,-1) {\small\(-1\)};
            \node at (0,-1) {\small\(\bullet\)};
            \node[above left, xshift=-6pt] at (-0.5,0) {\small\(-\frac{1}{2}\)};
            \node at (-0.5,0) {\small\(\bullet\)};
        \end{tikzpicture}
        \\
        (d)
        &
        \begin{tikzpicture}[scale=0.5]
            % Eixos
            \draw[->] (-3.0,0) -- (3.0,0) node[right] {\(x\)};
            \draw[->] (0,-3.0) -- (0,3.0) node[above] {\(y\)};

            % Inclinação errada
            \draw[domain=-2:2, thick, blue] plot (\x, {-\x + 1});

            \node[left] at (0,1) {\small\(1\)};
            \node at (0,1) {\small\(\bullet\)};
            \node[above right] at (1,0) {\small\(1\)};
            \node at (1,0) {\small\(\bullet\)};
        \end{tikzpicture}
        &
        (e) 
        &
        \begin{tikzpicture}[scale=0.5]
            % Eixos
            \draw[->] (-3.0,0) -- (3.0,0) node[right] {\(x\)};
            \draw[->] (0,-3.0) -- (0,3.0) node[above] {\(y\)};

            % Sinal invertido
            \draw[domain=-1.5:2.5, thick, blue] plot (\x, {\x - 1});

            \node[right] at (0,-1) {\small\(-1\)};
            \node at (0,-1) {\small\(\bullet\)};
            \node[above left] at (1,0) {\small\(1\)};
            \node at (1,0) {\small\(\bullet\)};
        \end{tikzpicture}
        &&
    \end{tabular}
\end{center}
}

% Define a resposta correta
\newcommand{\resposta}{A}

% Lógica condicional para exibição
\ifdefstring{\atividade}{lista}{%
    \alternativas
    \vspace{0.5em}
    
    \noindent\textbf{Resposta:} letra \textbf{\resposta}.
}{%
    \ifdefstring{\modo}{objetivo}{%
        \alternativas
    }{%
        % modo = discursiva → não mostra alternativas
    }
}
\end{exercicioBanco}


    \vspace{15pt}
    
    \phantomsection
    \addcontentsline{toc}{section}{Exercício 3}
    %%%%%%%%%%%%%%%%%%%%%%%%%%%%%%%%%
%
%       EXERCÍCIO
%
%%%%%%%%%%%%%%%%%%%%%%%%%%%%%%%%%

\ifdefstring{\atividade}{prova}{%
    \ifdefstring{\modo}{objetivo}{%
        \renewcommand{\valorquestao}{\ValorQObj\ ponto}
    }{%
        \renewcommand{\valorquestao}{\ValorQDisc\ pontos}
    }%
}%

\begin{exercicioBanco}[\valorquestao]
Dois reservatórios de água, inicialmente com o mesmo volume de $20$ litros, começam a ser abastecidos. Os volumes de água, $V_{1}$ e $V_{2}$, em função do tempo $t$, em minutos, estão representados no plano cartesiano pelas retas com interseção no ponto $A = (0,20)$ e que passam, respectivamente, pelos pontos $B = (10, 50)$ e $C = (30, 35)$.

\begin{center}
    \begin{tikzpicture}[scale=0.5]
        % Eixos
        \draw[->] (0,0) -- (6,0) node[below right] {\small{$t$ (min)}};
        \draw[->] (0,0) -- (0,6) node[above left] {\small{$V$ (L)}};
        
        % Funções
        \draw[domain=0:1.5, smooth, variable=\x, blue, thick] plot ({\x}, {3*\x+2}) node[below right] {\small $V_{1}$};
        \draw[domain=0:5, smooth, variable=\x, red, thick] plot ({\x}, {0.5*\x+2}) node[below right] {\small $V_{2}$};
        
        % Pontos
        \filldraw (0,2) circle (2pt) node[left] {\small $A$};
        \filldraw (1,5) circle (2pt) node[below right] {\small $B$};
        \filldraw (3,3.5) circle (2pt) node[below right] {\small $C$};
    \end{tikzpicture}
\end{center}

Sabe-se que, em determinado intervalo de tempo, o volume $V_{2}$ aumentou $5$ litros.

Nesse mesmo intervalo de tempo, determine quanto aumentou o volume $V_{1}$, em litros.


% Define as alternativas
\newcommand{\alternativas}{%
    \begin{center}
        \begin{tabularx}{\textwidth}{XXXXX}
            (a) \(10\). &
            (b) \(15\). &
            (c) \(20\). &
            (d) \(25\). &
            (e) \(30\).
        \end{tabularx}
    \end{center}
}

% Resposta correta
\newcommand{\resposta}{E}

% Lógica condicional para exibição
\ifdefstring{\atividade}{lista}{%
    \alternativas
    \vspace{0.5em}
    
    \noindent\textbf{Resposta:} letra \textbf{\resposta}.
}{%
    \ifdefstring{\modo}{objetivo}{%
        \alternativas
    }{%
        % modo = discursiva → não mostra alternativas
    }
}
\end{exercicioBanco}


    \vspace{15pt}
    
    \phantomsection
    \addcontentsline{toc}{section}{Exercício 4}
    %%%%%%%%%%%%%%%%%%%%%%%%%%%%%%%%%
%
%       EXERCÍCIO
%
%%%%%%%%%%%%%%%%%%%%%%%%%%%%%%%%%

\ifdefstring{\atividade}{prova}{%
    \ifdefstring{\modo}{objetivo}{%
        \renewcommand{\valorquestao}{\ValorQObj\ ponto}
    }{%
        \renewcommand{\valorquestao}{\ValorQDisc\ pontos}
    }%
}%

\begin{exercicioBanco}[\valorquestao]
Para organizar um grande estacionamento em um parque, um administrador precisa estimar a quantidade de funcionários necessários para controle de entrada e saída. Para cada 250 veículos estacionados, é necessário um funcionário.

Às 8h da manhã, o administrador observa que a área ocupada pelos veículos forma um retângulo de dimensões \(300\) m por \(200\) m. Sabe-se que, em média, cada veículo ocupa uma área de \(12\ \text{m}^2\).

Nas horas seguintes, espera-se que o número de veículos aumente a uma taxa de \(5.000\) veículos por hora até as 12h, quando o estacionamento será fechado.

Determine quantos funcionários serão necessários nesse horário.

% Define as alternativas
\newcommand{\alternativas}{%
    \begin{center}
        \begin{tabularx}{\textwidth}{XXXXX}
            (a) \(80\). &
            (b) \(100\). &
            (c) \(120\). &
            (d) \(140\). &
            (e) \(160\).
        \end{tabularx}
    \end{center}
}

% Resposta correta
\newcommand{\resposta}{B}

% Lógica condicional para exibição
\ifdefstring{\atividade}{lista}{%
    \alternativas
    \vspace{0.5em}
    
    \noindent\textbf{Resposta:} letra \textbf{\resposta}.
}{%
    \ifdefstring{\modo}{objetivo}{%
        \alternativas
    }{%
        % modo = discursiva → não mostra alternativas
    }
}
\end{exercicioBanco}


    \vspace{15pt}
    
    \phantomsection
    \addcontentsline{toc}{section}{Exercício 5}
    %%%%%%%%%%%%%%%%%%%%%%%%%%%%%%%%%
%
%       EXERCÍCIO
%
%%%%%%%%%%%%%%%%%%%%%%%%%%%%%%%%%

\ifdefstring{\atividade}{prova}{%
    \ifdefstring{\modo}{objetivo}{%
        \renewcommand{\valorquestao}{\ValorQObj\ ponto}
    }{%
        \renewcommand{\valorquestao}{\ValorQDisc\ pontos}
    }%
}%

\begin{exercicioBanco}[\valorquestao]
Considere as afirmações a seguir.
\begin{enumerate}[I.]
    \item Considere a função afim \(f(x) = ax\), com \(a \neq 0\). Então \(f(x) = -f(-x)\), para todo \(x \in \bbR\).
    %
    \item Seja \(f: \bbR \to \bbR\) uma função afim tal que a taxa de variação média no intervalo \([1,3]\) é igual a \(2\) e \(f(1) = 1\). Então \(f(4) = 7\).
    %
    \item Sejam \(f\) e \(g\) funções afins crescente e decrescente, respectivamente. Suponha que \(f(1) = g(1) = 0\). Então \(f(2)g(3) < 0\).
\end{enumerate}
%
Assinale a alternativa correta.

% Define as alternativas
\newcommand{\alternativas}{%
    \begin{center}
        \begin{tabularx}{\textwidth}{XXX}
            (a) Apenas I é verdadeira. &
            (b) I e II são verdadeiras. &
            (c) Apenas III é verdadeira. \\[5pt]
            (d) I e III são verdadeiras. &
            (e) Todas são verdadeiras.
        \end{tabularx}
    \end{center}
}

% Resposta correta
\newcommand{\resposta}{D}

% Lógica condicional para exibição
\ifdefstring{\atividade}{lista}{%
    \alternativas
    \vspace{0.5em}
    
    \noindent\textbf{Resposta:} letra \textbf{\resposta}.
}{%
    \ifdefstring{\modo}{objetivo}{%
        \alternativas
    }{%
        % modo = discursiva → não mostra alternativas
    }
}
\end{exercicioBanco}


    \vspace{15pt}
    
    \phantomsection
    \addcontentsline{toc}{section}{Exercício 6}
    %%%%%%%%%%%%%%%%%%%%%%%%%%%%%%%%%
%
%       EXERCÍCIO
%
%%%%%%%%%%%%%%%%%%%%%%%%%%%%%%%%%

\ifdefstring{\atividade}{prova}{%
    \ifdefstring{\modo}{objetivo}{%
        \renewcommand{\valorquestao}{\ValorQObj\ ponto}
    }{%
        \renewcommand{\valorquestao}{\ValorQDisc\ pontos}
    }%
}%

\begin{exercicioBanco}[\valorquestao]
Em uma cidade, dois aplicativos de transporte, $A$ e $B$, utilizam modelos distintos para calcular o preço de uma corrida em função da distância percorrida.

O aplicativo $A$ cobra uma taxa fixa de R\$ 6,00 mais R\$ 2,00 por quilômetro rodado. Já o aplicativo $B$ cobra uma taxa fixa de R\$ 10,00 mais R\$ 1,50 por quilômetro rodado.

Seja $x$ a distância, em quilômetros, de uma corrida. Determine para quais valores de $x$, em quilômetros, o aplicativo $A$ é mais vantajoso que o aplicativo $B$.

% Define as alternativas
\newcommand{\alternativas}{%
    \begin{center}
        \begin{tabularx}{\textwidth}{XXX}
            (a) \([8,\infty[\). &
            (b) \(]8,\infty[\). &
            (c) \([0,8]\). \\[5pt]
            (d) \([0,8[\). &
            (e) \([0, \infty[\).
        \end{tabularx}
    \end{center}
}

% Resposta correta
\newcommand{\resposta}{D}

% Lógica condicional para exibição
\ifdefstring{\atividade}{lista}{%
    \alternativas
    \vspace{0.5em}
    
    \noindent\textbf{Resposta:} letra \textbf{\resposta}.
}{%
    \ifdefstring{\modo}{objetivo}{%
        \alternativas
    }{%
        % modo = discursiva → não mostra alternativas
    }
}
\end{exercicioBanco}


    \vspace{15pt}
    
    %\newpage
    %\phantomsection
    %\addcontentsline{toc}{section}{Rascunho}
    %\input{../../Rascunho}
    
\end{document}
%*************** FINAL DO DOCUMENTO ***************
